% Consolidation322 draft. Sources319A and321C; not an assessed PDF.
\section{Growing marginal order at the critical moment budget}
\label{sec:growing-marginals}

The preceding marginal and moment constructions need not describe the same
weights. We now impose both requirements on one weighting law. Throughout
this section, $u_L$ is uniform probability on $\{-1,1\}^L$,
$0<\rho<1$ is fixed, and
\[
 s=\frac{1+\rho}{1-\rho},\qquad B_*=2^{1/s},\qquad
 Q=P_{\sqrt\rho},\qquad
 S_r=QM_rQ,\quad T_r=M_{\sqrt r}P_\rho M_{\sqrt r}.
\]
Here $P_t\chi_A=t^{|A|}\chi_A$. For $r\geq0$ with $\mathbb E r=1$,
$S_r$ and $T_r$ have the same eigenvalues, including zero multiplicities:
they are $AA^*$ and $A^*A$ for the square matrix $A=QM_{\sqrt r}$.
Also $\langle1,S_r1\rangle=1$, so $\lambda_1(T_r)\geq1$.
Exact uniform marginals of the weighting law $r u_L$ through order $k$
mean $\mathbb E[r\chi_A]=0$ for $1\leq|A|\leq k$.
This classical Fourier characterization concerns the weighting law;
no Perron stationary marginal assertion is made in this section.

Two different regimes result. At every fixed power $L^\alpha$ with
$\alpha<1/2$, exact matching is compatible with collapse at $B_*$.
At every fixed positive linear fraction of $L$, a bounded critical moment
instead forces convergence to the unweighted noise spectrum. The proof of
the latter statement uses the homogeneous moment inequality of
Kirshner and Samorodnitsky~\cite{kirshner-samorodnitsky}, not a new
hypercontractive inequality. The intervening sublinear regime is open here.

\begin{theorem}[Collapse with growing exact marginals]
\label{thm:growing-collapse}
Fix $0<\alpha<1/2$ and $0<b<1$. There are mean-one densities $r_L$,
for every integer $L\geq1$, such that
\[
 r_L\geq b,\qquad \|r_L\|_s\leq B_*,\qquad
 \mathbb E[r_L\chi_A]=0\quad(1\leq |A|\leq\lfloor L^\alpha\rfloor),
\]
while $\lambda_1(T_{r_L})\to1$, $\lambda_2(T_{r_L})\to1$ and
$\gap(T_{r_L})\to0$. Thus growing matching of these orders does not
raise the worst-case critical norm boundary $B_*$.
\end{theorem}

\begin{proof}
Choose $a=3/8-\alpha/4$, so $1/4<a<(1-\alpha)/2$, and fix a
sufficiently small $d\in(0,1)$. Set $\delta=dL^{-a}$ and
\[
 n(\delta)=\left(\frac{(1+\delta)^s+(1-\delta)^s}{2}\right)^{1/s},
 \quad R_L(x)=\prod_{i=1}^L\frac{1+\delta x_i}{n(\delta)},
 \quad m_L=\mathbb E R_L=n(\delta)^{-L}.
\]
The one-coordinate matrix and its largest eigenvalue are
\[
 H(\delta)=\frac1{n(\delta)}
 \begin{pmatrix}1&\delta\sqrt\rho\\\delta\sqrt\rho&\rho\end{pmatrix},
 \qquad
 \kappa(\delta)=\frac{1+\rho+
 \sqrt{(1-\rho)^2+4\rho\delta^2}}{2n(\delta)}.
\]
Writing $D=\diag(1,\rho)$ and
$\tau(\delta)=\|D^{1/2}H(\delta)^{1/2}\|$, direct diagonalization gives
\[
 \tau(\delta)^2=\frac{1+\rho^2+
 \sqrt{(1-\rho^2)^2+4\rho^2\delta^2}}{2n(\delta)}.
\]
At the fixed critical exponent the Taylor expansions are
\[
 \log n(\delta)=\frac{s-1}{2}\delta^2+O(\delta^4),\quad
 \log\kappa(\delta)=O(\delta^4),\quad
 \log\tau(\delta)=-\frac{\rho\delta^2}{2(1-\rho^2)}+O(\delta^4).
\]
Here $L\delta^4=d^4L^{\alpha-1/2}\to0$, whereas
$L\delta^2=d^2L^{1/4+\alpha/2}\to\infty$.
Consequently $c_L=\kappa(\delta)^{-L}\to1$, $m_L\to0$, and
$Z_L=1+c_Lm_L\to1$.

For $v_L=(1+c_LR_L)/Z_L$, the unnormalized operator is
$F_L=P_\rho+c_LH(\delta)^{\otimes L}$. Both summands have top
eigenvalue one. Form the Gram matrix of their positive square roots.
Its diagonal blocks are these two summands and its cross-block norm is
$\eta_L=\sqrt{c_L}\tau(\delta)^L\to0$. The block diagonal matrix has
its first two eigenvalues exactly one. Operator-norm perturbation puts
the first two Gram eigenvalues in $[1-\eta_L,1+\eta_L]$.
For large $L$ these are positive; the nonzero Gram spectrum is the spectrum
of $F_L$. Division by $Z_L$ therefore proves that the actual first two
eigenvalues of $S_{v_L}$ tend to one.

We also have $\mathbb E R_L^s=1$. For each fixed $0<t<s$,
\[
 \log\mathbb E R_L^t=-\frac{t(s-t)}2L\delta^2+O(L\delta^4)
 \longrightarrow-\infty.
\]
With $\theta=\min(1,s/2)$, the mean-value theorem for $y\leq1$ and
reciprocal substitution for $y\geq1$ give
\[
 0\leq(1+y)^s-1-y^s\leq
 s2^{s-1}(y^\theta+y^{s-\theta})\quad(y\geq0).
\]
Both exponents lie strictly between zero and $s$. Applying this inequality
to $c_LR_L$ yields $\mathbb E(1+c_LR_L)^s\to2$, hence
$\|v_L\|_s\to B_*$.

Put $k_L=\lfloor L^\alpha\rfloor$ and erase the nonconstant low levels:
\[
 p_L=m_L\sum_{1\leq|A|\leq k_L}\delta^{|A|}\chi_A,\qquad
 w_L=\frac{1+c_L(R_L-p_L)}{Z_L}.
\]
The exact expansion $R_L=m_L\sum_A\delta^{|A|}\chi_A$ proves that
$w_L$ has mean one and the required zero coefficients. Moreover
\[
 \|p_L\|_\infty\leq m_L\sum_{j=1}^{k_L}\binom Lj\delta^j
 \leq m_L k_LL^{k_L}\longrightarrow0.
\]
Indeed $\log m_L=-(s-1)d^2L^{1/4+\alpha/2}/2+O(L\delta^4)$,
and $\log k_L+k_L\log L=o(L^{1/4+\alpha/2})$.
It follows that $w_L\geq(1-c_L\|p_L\|_\infty)/Z_L\to1$ in this
lower-bound sense, so eventually $w_L\geq b$. Also
$\|w_L-v_L\|_\infty\to0$, $\|w_L\|_s\to B_*$, and
$\|S_{w_L}-S_{v_L}\|\leq\|w_L-v_L\|_\infty\to0$.
The uniform baseline, together with this sup-norm bound, supplies the
positivity that an unrestricted Fourier erasure would not preserve.

Finally let $B_L=\|w_L\|_s$, and set $\varepsilon_L=0$ when
$B_L\leq B_*$, otherwise
$\varepsilon_L=(B_L-B_*)/(B_L-1)$. On the latter event the denominator
is positive, $0<\varepsilon_L<1$, and $\varepsilon_L\to0$.
Then $r_L=(1-\varepsilon_L)w_L+\varepsilon_L$ retains the exact
marginals, mean, and eventual floor. Minkowski gives $\|r_L\|_s\leq B_*$.
Endpoint hypercontractivity and H\"older, with
$p=2s/(s-1)=1+1/\rho$, imply for every real signed $f$ that
\begin{equation}\label{eq:critical-signed-322}
 \|QM_fQ\|\leq\|f\|_s,
 \qquad \|Q\|_{2\to p}\leq1.
\end{equation}
Thus $\|S_{r_L}-S_{w_L}\|\leq\varepsilon_L(B_L+1)\to0$.
The actual leading eigenvalues still tend to one. Replace the finitely
many densities before all positivity and floor conditions hold by $r_L=1$.
Every stated constraint then holds in every dimension. The positive kernel
has a simple Perron root in each dimension, while its normalized gap tends
to zero. Below $B_*$ the earlier sign-compression bound
$2\lambda_2^s\leq\mathbb E r^s$ already supplies a uniform gap.
\end{proof}

\begin{theorem}[Rigidity from critical-power uniform integrability]
\label{thm:marginal-ui}
Suppose $k_L\to\infty$, $r_L\geq0$, $\mathbb E r_L=1$, and the
weighting laws have uniform marginals through order $k_L$. If
\[
 \lim_{M\to\infty}\sup_L
 \mathbb E[r_L^s\mathbf1_{r_L>M}]=0,
\]
then $\|S_{r_L}-P_\rho\|\to0$ and $\gap(T_{r_L})\to1-\rho$.
For $\rho<\sigma<1$, put $t=\sqrt{\rho/\sigma}$,
$m_L=\lfloor k_L/2\rfloor$, and
$\eta_M=\sup_L\|r_L\mathbf1_{r_L>M}\|_s$. More precisely,
\[
 \|S_{r_L}-P_\rho\|
 \leq2(M+1)t^{m_L+1}+2\eta_M.
\]
\end{theorem}
\begin{proof}
Let $\Pi$ project onto Walsh degrees at most $m_L$. With $h=r_L-1$,
the $(A,C)$ entry of $QM_hQ$ is
$\rho^{(|A|+|C|)/2}\widehat h(A\mathbin\triangle C)$.
Since $|A\mathbin\triangle C|\leq2m_L\leq k_L$ on this block,
including the zero constant coefficient, $\Pi QM_hQ\Pi=0$.

Split $h=h_b+h_t$, with $h_b=r_L\mathbf1_{r_L\leq M}-1$ and
$h_t=r_L\mathbf1_{r_L>M}$. Write $Q_\sigma=P_{\sqrt\sigma}$ and
$D_t=P_t$. For $B_b=QM_{h_b}Q$ we have
$B_b=D_t A_bD_t$, $A_b=Q_\sigma M_{h_b}Q_\sigma$,
$\|A_b\|\leq M+1$, and $\|D_t(I-\Pi)\|\leq t^{m_L+1}$.
The identity
$B_b-\Pi B_b\Pi=(I-\Pi)B_b+\Pi B_b(I-\Pi)$ bounds its norm by
$2(M+1)t^{m_L+1}$. For $B_t=QM_{h_t}Q$,
\eqref{eq:critical-signed-322} gives $\|B_t-\Pi B_t\Pi\|\leq2\eta_M$.
Since $\Pi(B_b+B_t)\Pi=0$, addition proves the estimate.
First send $L\to\infty$ at fixed $M$, then $M\to\infty$.
Self-adjoint perturbation and the actual spectrum $1,\rho,\ldots$ of
$P_\rho$ give the first two eigenvalue limits $1,\rho$ and the gap limit.
\end{proof}

For comparison, a strictly higher moment gives an explicit rate without
truncation. If $q>s$, $\|r_L\|_q\leq B$, and
$\rho<\sigma\leq(q-1)/(q+1)$, the same factorization gives
\[
 \|S_{r_L}-P_\rho\|\leq2(B+1)
 (\rho/\sigma)^{(\lfloor k_L/2\rfloor+1)/2}.
\]
Here $Q_\sigma M_{r_L-1}Q_\sigma$ has norm at most $B+1$ by
H\"older and hypercontractivity. The corresponding low block vanishes
exactly. This estimate and the preceding theorem are sufficient conditions,
not characterizations of all stable families.

\begin{theorem}[Rigidity for every positive linear marginal fraction]
\label{thm:linear-marginal-rigidity}
Suppose $r_L\geq0$, $\mathbb E r_L=1$, $\|r_L\|_s\leq B<\infty$,
and every marginal through order $k_L\leq L$ is uniform. If
$\liminf_L k_L/L>0$, then $\|S_{r_L}-P_\rho\|\to0$.
In particular the two leading eigenvalues of $T_{r_L}$ tend to $1,\rho$,
and $\gap(T_{r_L})\to1-\rho$. No density floor, uniform integrability,
or supercritical moment assumption is required.
\end{theorem}

\begin{proof}
Set $p=2s/(s-1)=1+1/\rho>2$. We use Theorem~1.3 and
Proposition~2.7 of the version-pinned primary source
\cite{kirshner-samorodnitsky}. In uniform probability normalization, its
continuous function $\psi(p,x)$ on $[0,1/2]$ bounds every real homogeneous
degree-$j$ Walsh polynomial, $j\leq L/2$, by
\[
 \|f_j\|_p\leq2^{L\psi(p,j/L)/p}\|f_j\|_2.
\]
For fixed $p$, $\psi(p,0)=0$, the right derivative at zero is
$(p/2)\log_2(p-1)$, and the function is strictly concave.
These are attributed external mathematical inputs. Strict concavity gives
\[
 \Delta_p(x)=\frac{x\log(p-1)}2-
 \frac{\log2}{p}\psi(p,x)>0\quad(0<x\leq1/2).
\]
For completeness, if $f(0)=0$ is strictly concave with right derivative
$d$, then $f(x)/x<f(x/2)/(x/2)\leq d$; this also handles the upper
endpoint. Natural logarithms are used except for the displayed $\log_2$.

For $x\in[0,1]$ let $y=\min(x,1-x)$ and
\[
 a_p(x)=\frac{x\log(p-1)}2-\frac{\log2}{p}\psi(p,y)
 =\Delta_p(y)+\frac{(x-y)\log(p-1)}2.
\]
This is continuous and strictly positive for every $x>0$, including
$a_p(1)=\log(p-1)/2$. For each fixed $0<\gamma\leq1$, compactness gives
$a_\gamma=\min_{x\in[\gamma,1]}a_p(x)>0$.
For $j\leq L/2$ the published inequality and $\rho=1/(p-1)$ imply
\[
 \|Qf_j\|_p\leq e^{-La_p(j/L)}\|f_j\|_2.
\]
For $j>L/2$, multiplication by the full parity character maps the entire
degree-$j$ level to degree $L-j$ while preserving all $L^q$ norms.
Applying the preceding homogeneous inequality there and retaining the
original factor $\rho^{j/2}$ gives the same bound. This does not assert
that parity multiplication commutes with $Q$.

Let $H_\gamma$ project onto degrees $j\geq\gamma L$. Triangle inequality,
then Cauchy--Schwarz over at most $L+1$ orthogonal levels, give
\begin{equation}\label{eq:high-level-endpoint-322}
 \|QH_\gamma\|_{2\to p}\leq\sqrt{L+1}\,e^{-La_\gamma}.
\end{equation}
Take fixed $0<c<\liminf k_L/L$ and put $\gamma=c/2$. For all sufficiently
large $L$, the complement $H=I-\Pi$ of degrees at most
$m_L=\lfloor k_L/2\rfloor$ has degrees
$j\geq m_L+1\geq k_L/2\geq\gamma L$.
For $\mathcal B=QM_{r_L-1}Q$, exact marginals give $\Pi\mathcal B\Pi=0$,
hence $\mathcal B=H\mathcal B+\Pi\mathcal BH$.
For real unit vectors $u,v$, H\"older with $1/s+2/p=1$ yields
\[
 |\langle u,H\mathcal Bv\rangle|
 \leq\|r_L-1\|_s\|QHu\|_p\|Qv\|_p
 \leq(B+1)\sqrt{L+1}\,e^{-La_\gamma}.
\]
The other term uses $Q\Pi u$ and $QHv$. Since $\Pi$ is an $L^2$
contraction, $\|Q\Pi\|_{2\to p}\leq1$; no $L^p$ contraction property
of $\Pi$ is needed. Thus
\[
 \|S_{r_L}-P_\rho\|\leq2(B+1)\sqrt{L+1}\,e^{-La_\gamma}.
\]
This tends to zero. The ordered-eigenvalue perturbation bound applies in
the same finite dimensional space for each $L$. It gives the actual
top-two limits, including when $r_L$ has zeros, and the common spectrum
of $S_{r_L},T_{r_L}$ supplies the asserted gap limit.
\end{proof}

\paragraph{Scope of the separation.}
Theorems~\ref{thm:growing-collapse} and~\ref{thm:linear-marginal-rigidity}
operate at the same critical cap $B_*$. Their conclusions do not locate a
sharp growth threshold between sub-square-root powers and linear order.
In particular qualitative strict concavity does not provide a uniform
quadratic deficit near zero, and the constant $a_\gamma$ above is fixed
only for fixed $\gamma>0$. No shrinking-cutoff rate is inferred.
Theorem~\ref{thm:marginal-ui} supplies a separate sufficient condition at
arbitrary diverging marginal order. All statements concern finite Boolean
transfer operators with global weights. They imply no fixed-range Gibbs,
local-dynamics, physical-gap, continuum, or Yang--Mills result.
The classical marginal/Fourier criterion and coefficient erasure are prior
methods; the homogeneous moment improvement is explicitly attributed.
Historical priority for their joint spectral application remains unestablished.
